\section{Fundamentação Teórica}

\subsection{Modelos de linguagem e o objetivo de pré-treino}

Um modelo de linguagem autorregressivo, baseado na arquitetura
\textit{Transformer}~\cite{vaswani2017attention}, estima a probabilidade do próximo
token $x_t$ dado o contexto anterior, $P(x_t \mid x_{<t})$. O treinamento minimiza a
entropia cruzada média entre a distribuição prevista e o token observado:

\begin{equation}
  \mathcal{L} = -\frac{1}{N} \sum_{t=1}^{N} \log P(x_t \mid x_{<t}),
  \label{eq:entropia}
\end{equation}

\noindent onde $N$ é o número de tokens. O \texttt{Qwen2.5-0.5B} adotado neste
trabalho possui 494 milhões de parâmetros, 24 camadas, dimensão oculta de 896 e
vocabulário de 151.936 tokens~\cite{qwen2025}.

\subsection{Pré-treino continuado (\textit{continued pretraining})}

O pré-treino continuado consiste em prosseguir o treino não supervisionado de um
modelo já pré-treinado sobre um corpus de um domínio-alvo, usando o mesmo objetivo
da Equação~\ref{eq:entropia}. A técnica adapta o modelo ao vocabulário, ao estilo e
às regularidades estatísticas do domínio sem necessidade de rótulos, e é
reconhecidamente eficaz para domínios especializados~\cite{gururangan2020dont}. É a
abordagem da Questão 01.

\subsection{Pós-treino: ajuste fino supervisionado (SFT)}

O \textit{Supervised Fine-Tuning} (SFT) treina o modelo sobre pares de instrução e
resposta, ensinando-o a seguir comandos e a responder no formato
desejado~\cite{ouyang2022instructgpt}. Os exemplos seguem tipicamente um molde de
prompt como o do \textit{Stanford Alpaca}~\cite{taori2023alpaca}, com campos de
instrução, entrada opcional e resposta. No SFT clássico (Questão 02), todos os
parâmetros do modelo são atualizados, o que oferece máxima capacidade de adaptação
ao custo de elevada demanda de memória e tempo.

\subsection{Adaptação eficiente: LoRA e QLoRA}

A \textit{Low-Rank Adaptation} (LoRA)~\cite{hu2022lora} congela os pesos originais
$W_0$ e aprende apenas um par de matrizes de baixo posto $A$ e $B$, de modo que a
atualização efetiva seja $W_0 + BA$, com $\mathrm{posto}(BA) = r \ll d$. Como apenas
$A$ e $B$ são treinados, o número de parâmetros atualizados cai em ordens de
grandeza, reduzindo memória e tempo sem alterar a arquitetura na inferência.

O QLoRA~\cite{dettmers2023qlora} estende a ideia carregando o modelo base
\emph{quantizado} em 4 bits (formato NF4) e aplicando os adaptadores LoRA sobre essa
representação comprimida~\cite{dettmers2022llmint8}. Isso permite ajustar modelos em
GPUs modestas. Ambas as técnicas são disponibilizadas pela biblioteca
PEFT~\cite{peft2022} e correspondem à Questão 03.

\subsection{Métricas de avaliação}
\label{sec:metricas}

As três questões usam o mesmo conjunto de métricas, calculadas sobre um conjunto de
avaliação reservado (\textit{held-out}):

\begin{itemize}[leftmargin=2.2em]
  \item \textbf{Entropia cruzada:} a perda média da Equação~\ref{eq:entropia};
        quanto menor, melhor o ajuste do modelo ao texto.
  \item \textbf{Perplexidade:} o exponencial da entropia cruzada,
        $\mathrm{PPL} = e^{\mathcal{L}}$~\cite{jelinek1977perplexity}. Interpreta-se
        como o número médio de opções entre as quais o modelo hesita por token; uma
        perplexidade próxima de 1 indica previsão quase determinística.
  \item \textbf{Acurácia de previsão de tokens:} a fração de tokens em que o token de
        maior probabilidade (\textit{argmax}) coincide com o token real.
\end{itemize}
